\documentclass{amsart}
% For dealing with references we use the comment environment
\usepackage{verbatim}
\newenvironment{reference}{\comment}{\endcomment}
%\newenvironment{reference}{}{}
\newenvironment{slogan}{\comment}{\endcomment}
\newenvironment{history}{\comment}{\endcomment}

% For commutative diagrams we use Xy-pic
\usepackage[all]{xy}

% We use 2cell for 2-commutative diagrams.
\xyoption{2cell}
\UseAllTwocells

% We use multicol for the list of chapters between chapters
\usepackage{multicol}

% This is generally recommended for better output
\usepackage{lmodern}
\usepackage[T1]{fontenc}

% For cross-file-references

% Package for hypertext links:
\usepackage{hyperref}

% For any local file, say "hello.tex" you want to link to please

% Theorem environments.
%
\theoremstyle{plain}
\newtheorem{theorem}[subsection]{Theorem}
\newtheorem{proposition}[subsection]{Proposition}
\newtheorem{lemma}[subsection]{Lemma}

\theoremstyle{definition}
\newtheorem{definition}[subsection]{Definition}
\newtheorem{example}[subsection]{Example}
\newtheorem{exercise}[subsection]{Exercise}
\newtheorem{situation}[subsection]{Situation}

\theoremstyle{remark}
\newtheorem{remark}[subsection]{Remark}
\newtheorem{remarks}[subsection]{Remarks}

\numberwithin{equation}{subsection}

% Macros
%
\def\lim{\mathop{\mathrm{lim}}\nolimits}
\def\colim{\mathop{\mathrm{colim}}\nolimits}
\def\Spec{\mathop{\mathrm{Spec}}}
\def\Hom{\mathop{\mathrm{Hom}}\nolimits}
\def\Ext{\mathop{\mathrm{Ext}}\nolimits}
\def\SheafHom{\mathop{\mathcal{H}\!\mathit{om}}\nolimits}
\def\SheafExt{\mathop{\mathcal{E}\!\mathit{xt}}\nolimits}
\def\Sch{\mathit{Sch}}
\def\Mor{\mathop{\mathrm{Mor}}\nolimits}
\def\Ob{\mathop{\mathrm{Ob}}\nolimits}
\def\Sh{\mathop{\mathit{Sh}}\nolimits}
\def\NL{\mathop{N\!L}\nolimits}
\def\CH{\mathop{\mathrm{CH}}\nolimits}
\def\proetale{{pro\text{-}\acute{e}tale}}
\def\etale{{\acute{e}tale}}
\def\QCoh{\mathit{QCoh}}
\def\Ker{\mathop{\mathrm{Ker}}}
\def\Im{\mathop{\mathrm{Im}}}
\def\Coker{\mathop{\mathrm{Coker}}}
\def\Coim{\mathop{\mathrm{Coim}}}

% Boxtimes
%
\DeclareMathSymbol{\boxtimes}{\mathbin}{AMSa}{"02}

%
% Macros for moduli stacks/spaces
%
\def\QCohstack{\mathcal{QC}\!\mathit{oh}}
\def\Cohstack{\mathcal{C}\!\mathit{oh}}
\def\Spacesstack{\mathcal{S}\!\mathit{paces}}
\def\Quotfunctor{\mathrm{Quot}}
\def\Hilbfunctor{\mathrm{Hilb}}
\def\Curvesstack{\mathcal{C}\!\mathit{urves}}
\def\Polarizedstack{\mathcal{P}\!\mathit{olarized}}
\def\Complexesstack{\mathcal{C}\!\mathit{omplexes}}
% \Pic is the operator that assigns to X its picard group, usage \Pic(X)
% \Picardstack_{X/B} denotes the Picard stack of X over B
% \Picardfunctor_{X/B} denotes the Picard functor of X over B
\def\Pic{\mathop{\mathrm{Pic}}\nolimits}
\def\Picardstack{\mathcal{P}\!\mathit{ic}}
\def\Picardfunctor{\mathrm{Pic}}
\def\Deformationcategory{\mathcal{D}\!\mathit{ef}}

\setlength{\emergencystretch}{3em}
\begin{document}
\title[Stacks Project, Tag 0FD0]{366 --- Extension of coherent sheaf morphisms after modifying their domain off an open subset}
\maketitle
\noindent Copyright (C) 2005--2025 Johan de Jong. Stacks Project authors. Source: \href{https://stacks.math.columbia.edu/tag/0FD0}{Tag 0FD0}.

\noindent Editorial difficulty note (not author text). Difficulty H1: The proof must globalize extension from a Noetherian local region to a merely locally Noetherian scheme. It constructs a partially ordered system of compatible coherent domain subsheaves and extended maps, proves gluing along chains, and uses local extension to rule out a proper maximal domain..

\noindent Editorial provenance note (not author text). History: excerpt from revision \texttt{a0444\allowbreak 6e57e\allowbreak c1fbc\allowbreak 252a8\allowbreak 71afc\allowbreak ec775\allowbreak 2fb28\allowbreak 07b14}, coherent.tex, lines 2745--2791. Reference presentation was adapted; original-author context and core support blocks were retained or added. The mathematical wording is unchanged. Licensed under GNU FDL 1.2 or later; no invariant sections or cover texts. See ../licenses/COPYING. Context blocks reproduce author definitions and supporting results; they are not additional counted problems.

\begin{lemma}
\label{lemma-extend-coherent}
Let $X$ be a locally Noetherian scheme. Let $\mathcal{F}$, $\mathcal{G}$
be coherent $\mathcal{O}_X$-modules. Let $U \subset X$ be open and
let $\varphi : \mathcal{F}|_U \to \mathcal{G}|_U$ be an
$\mathcal{O}_U$-module map. Then there exists a coherent
submodule $\mathcal{F}' \subset \mathcal{F}$ agreeing with
$\mathcal{F}$ over $U$ such that $\varphi$ extends to
$\varphi' : \mathcal{F}' \to \mathcal{G}$.
\end{lemma}

\begin{proof}
Let $\mathcal{I} \subset \mathcal{O}_X$ be the coherent sheaf of ideals
cutting out the reduced induced scheme structure on $X \setminus U$.
If $X$ is Noetherian, then Lemma \href{https://stacks.math.columbia.edu/tag/01YB}{lemma homs over open (Tag 01YB)} tells us that
we can take $\mathcal{F}' = \mathcal{I}^n\mathcal{F}$
for some $n$. The general case will follow from this using Zorn's lemma.

\medskip\noindent
Consider the set of triples $(U', \mathcal{F}', \varphi')$ where
$U \subset U' \subset X$ is open, $\mathcal{F}' \subset \mathcal{F}|_{U'}$
is a coherent subsheaf agreeing with $\mathcal{F}$ over $U$, and
$\varphi' : \mathcal{F}' \to \mathcal{G}|_{U'}$ restricts to
$\varphi$ over $U$. We say
$(U'', \mathcal{F}'', \varphi'') \geq (U', \mathcal{F}', \varphi')$
if and only if $U'' \supset U'$, $\mathcal{F}''|_{U'} = \mathcal{F}'$,
and $\varphi''|_{U'} = \varphi'$.
It is clear that if we have a totally ordered collection
of triples $(U_i, \mathcal{F}_i, \varphi_i)$, then we
can glue the $\mathcal{F}_i$ to a subsheaf $\mathcal{F}'$ of $\mathcal{F}$
over $U' = \bigcup U_i$ and extend $\varphi$ to a map
$\varphi' : \mathcal{F}' \to \mathcal{G}|_{U'}$.
Hence any totally ordered subset of triples has an upper bound.
Finally, suppose that $(U', \mathcal{F}', \varphi')$
is any triple but $U' \not = X$. Then we can choose an
affine open $W \subset X$ which is not contained in $U'$.
By the result of the first paragraph we can extend
the subsheaf $\mathcal{F}'|_{W \cap U'}$ and the restriction
$\varphi'|_{W \cap U'}$ to some subsheaf $\mathcal{F}'' \subset \mathcal{F}|_W$
and map $\varphi'' : \mathcal{F}'' \to \mathcal{G}|_W$.
Of course the agreement between $(\mathcal{F}', \varphi')$ and
$(\mathcal{F}'', \varphi'')$ over $W \cap U'$ exactly means that
we can extend this to a triple $(U' \cup W, \mathcal{F}''', \varphi''')$.
Hence any maximal triple $(U', \mathcal{F}', \varphi')$
(which exist by Zorn's lemma) must have $U' = X$ and the
proof is complete.
\end{proof}
\end{document}
